\chapter{Metodología Aplicada}

% ── 5.1 ──────────────────────────────────────────────────
\section{Metodología de Desarrollo}

Para el desarrollo del proyecto se adoptó la metodología ágil \textbf{Scrum},
dado que permite una alta flexibilidad y adaptación a los cambios, además de
fomentar una comunicación constante entre los miembros del equipo. Scrum se
basa en ciclos de trabajo cortos y entregables funcionales (sprints) que
permiten un seguimiento continuo del progreso.

El proceso de desarrollo sigue las siguientes fases dentro del marco de Scrum:

\begin{itemize}[leftmargin=1.5em]
  \item \textbf{Product Backlog:} Lista priorizada de los requisitos del
        proyecto, mantenida y priorizada por el Product Owner.
  \item \textbf{Sprint Planning:} Antes de cada sprint el equipo selecciona
        los ítems del backlog a completar y se compromete a entregarlos.
  \item \textbf{Daily Stand-Up:} Reuniones diarias breves para revisar
        progreso, detectar bloqueos y ajustar tareas.
  \item \textbf{Sprint Review:} Al final del sprint se presenta el incremento
        al Product Owner y partes interesadas.
  \item \textbf{Sprint Retrospective:} El equipo identifica áreas de mejora
        en el proceso y aplica ajustes al siguiente sprint.
\end{itemize}

\section{Ciclo de Vida del Proyecto}

El proyecto se organizó en \textbf{3 sprints de 4 semanas} cada uno:

\begin{longtable}{>{\bfseries}p{1.5cm} p{11.5cm}}
  \toprule
  \textbf{Sprint} & \textbf{Alcance} \\
  \midrule
  \endhead

  1 & Configuración del proyecto, diseño del schema en Supabase,
      autenticación, RLS, CRUD del menú e integración con Storage
      para manejo de imágenes. \\[4pt]

  2 & Registro de pedidos, vista de mesas en tiempo real mediante
      Realtime, manejo de estados de pedido, módulo de gastos,
      funciones SQL de reportes y ranking de platos. \\[4pt]

  3 & Carta digital pública por QR, sesión anónima por mesa,
      llamada al mesero, confirmación de pagos por admin,
      pruebas finales y despliegue con EAS Build. \\

  \bottomrule
  \caption{Release Plan — Sprints del proyecto}
\end{longtable}

\section{Herramientas de Gestión}

\begin{itemize}[leftmargin=1.5em]
  \item \textbf{Notion:} Gestión del sprint backlog con vista Kanban
        (Por hacer / En progreso / En revisión / Completado),
        centralización de documentación técnica y seguimiento de tareas.
  \item \textbf{GitHub:} Control de versiones distribuido, revisión de código
        mediante pull requests y entrega de incrementos de forma controlada.
  \item \textbf{Discord:} Canal principal de comunicación cotidiana,
        dividido en espacios temáticos (desarrollo, diseño, pruebas).
  \item \textbf{Excalidraw:} Diagramación colaborativa de esquemas de base de
        datos, flujos de navegación y arquitectura del sistema.
  \item \textbf{Postman:} Testing y documentación de llamadas RPC y REST
        durante el desarrollo e integración.
\end{itemize}